\section{From contour convergence to planar topology}
\label{bg:topology-section}

This section turns uniform contour convergence into topological control
of the discrete map. We first join any prescribed lifts of nearby points
by paths with small projected images. These joins yield disk fillings,
winding separation, and actual intersections of primal crossings. We
then construct cellular neighborhoods with complete incidence layers
and preserve them when a bounded contour window is capped to a sphere.
No bound on graph length or normalized passage cost is asserted here.

We use the round distance $d_{\mathrm{sph}}$ on $\mathbb S^2$ and write
$B(x,r)$ for a round ball; its normalization is immaterial. All round
balls below have radius smaller than a fixed convexity radius. Graphs
on surfaces are embedded, so distinct open edges are disjoint; loops
and multiple edges are allowed.

\subsection{Contour links and small-image joins}

Recall the resolved tree/cotree sewing convention. A chronological mesh
time specifies a primal corner. Matched up- and down-steps of the first
coordinate identify primal corners in the contour tree. A matched pair
in the second coordinate bounds a sewn quadrangle whose primal diagonal
is an all-edge primal edge. Its endpoint corners occur at the two
steps, up to one chronological mesh step. Thus a limiting
second-coordinate identification must be approximated by a
\emph{matched discrete pair} before it supplies a primal edge.

\begin{lemma}[Approximating finitely many horizontal generators]
\label{bg:chord-approximation}
Suppose resolved discrete contour pairs converge uniformly to a
continuous contour pair on a compact chronological interval. The time
and height meshes tend to zero. Assume that, in each limiting scalar
contour, local minima are strict and have distinct heights. End
identifications are treated by the usual excursion convention; in
particular, the two ends of a positive excursion cause no additional
interior zero.

Fix a finite list of horizontal generators $(s,t)$, each satisfying
\[
 X(s)=X(t)=\min_{[s,t]}X
\]
for one of the two coordinates $X$. Each chosen pair is contained
in the interior of the interval of uniform convergence, so that both
exterior sides are available. The only endpoint exception is the
pair of genuine endpoints of a strictly positive excursion. For every generator there are
chronological endpoints $s_j\to s$, $t_j\to t$ such that their discrete
primal corners are identified, in the first-coordinate case, or are
joined by one genuine all-edge primal edge, in the second-coordinate
case. The choices can be made simultaneously for the finite list.
More generally, if a generator is specified as a fixed finite list of
such scalar generators, every member of that list has this property.

The conclusion concerns individual links. Approximations to a shared
intermediate time need not be the same discrete vertex.
\end{lemma}

\begin{proof}
Write $a=X(s)=X(t)$. First suppose $X>a$ on $(s,t)$.
For a level $a+\delta$ approaching $a$ from above, chosen to avoid
the countable set of local-minimum heights, take the component
of $\{X>a+\delta\}$ containing a fixed interior compact subinterval.
Its endpoints tend to $s,t$ as $\delta\downarrow0$. Indeed, $X-a$ has
a strictly positive minimum on each compact subinterval of $(s,t)$,
whereas its endpoint values are zero. At each boundary of the chosen
component there are strictly lower values arbitrarily close on the
exterior side: otherwise that boundary would be a local minimum at
the excluded level. Thus, for fixed sufficiently small $\delta$,
uniform contour convergence supplies a discrete excursion
above a nearby lattice level whose entrance and exit are as close to
these endpoints as desired. The entrance up-step and exit down-step
are a matched pair. Let $\delta$ tend to zero sufficiently slowly with
the discrete index. The sewing dictionary gives the claimed
identification or primal diagonal, with endpoint errors tending to
zero.

There is at most one time $r\in(s,t)$ with $X(r)=a$: two such times
would be local minima of the same height. If one exists, there are
values below $a$ arbitrarily close to $s$ on its exterior side and to
$t$ on its exterior side. Otherwise $s$ or $t$ would itself be a local
minimum of height $a$, distinct from $r$. For $\delta>0$, again
avoiding local-minimum heights, take instead
the component of $\{X>a-\delta\}$ containing $[s,t]$. The just noted
exterior values show that its endpoints tend to $s,t$ as
$\delta\downarrow0$. Apply the same discrete matched-step argument
at a nearby lattice level below $a$. This proves the assertion for
the outer pair at a scalar branch triple. Each consecutive pair at
that triple has strictly positive excursion interior and is covered
by the first argument. The excursion end convention is covered by
the first argument as well. A constant generator needs no edge.

For finitely many prescribed generators, take one diagonal choice of
the level errors and the discrete index thresholds. The number of
matched-step links per prescribed scalar generator is one and does
not depend on the discrete index. No rate of convergence is needed.
\end{proof}

For Brownian contours the scalar hypothesis holds almost surely.
Indeed, the minimum on a fixed interval is almost surely unique,
and minima on disjoint rational intervals have distinct values.
Every local minimum is the minimum on a rational interval, so a
countable intersection gives strictness and distinct heights.
Local absolute continuity proves the same assertion on compact
interior intervals of a quadrant excursion; exhaust its open time
interval and use the excursion convention at the ends.

To pass from finitely many contour links to uniform joins, we use the
following deterministic assumptions. Let $\Sigma_j$ be triangulated
topological spheres with embedded primal graphs $G_j\subset\Sigma_j$,
and let
\[
 \pi_j:\Sigma_j\longrightarrow\mathbb S^2,
 \qquad c_j:[0,1]\longrightarrow G_j,
 \qquad \eta:[0,1]\longrightarrow\mathbb S^2
\]
be continuous maps. The chronological walk $c_j$ visits every primal
vertex, and $\eta$ is onto. We require
\begin{enumerate}[label=(T\arabic*)]
\item $\sup_t d_{\mathrm{sph}}(\pi_j(c_j(t)),\eta(t))\to0$.
\item Every point $x\in\Sigma_j$ can be connected to some primal
vertex $v$ by a path whose projected diameter tends to zero uniformly
in $x$. For example, this follows if every closed triangle contains a
primal vertex and
$\max_{\sigma}\operatorname{diam}(\pi_j(\sigma))\to0$.
\item If $\eta(s)=\eta(t)$, then $s,t$ are connected by a finite chain
of limiting horizontal generators. For each generator $(a,b)$ in
such a chain there are $a_j\to a$, $b_j\to b$ and a primal path from
$c_j(a_j)$ to $c_j(b_j)$ whose projected diameter tends to zero.
\end{enumerate}
In (T3), the finite chain may depend on $s,t$; no uniform bound on
chain length is required.
Lemma~\ref{bg:chord-approximation}, together with vanishing projected
edge diameters, supplies the link approximation in (T3).

Here is the continuous projection used for the sewn maps. Assign each
primal or dual vertex the location $\eta(t)$ of one of its visits.
All visits to that vertex have uniformly vanishing disagreement:
otherwise a contrary sequence of pairs of visit times would converge
subsequentially, and uniform contour convergence would pass their
scalar equality or matching relation to a scalar chord with equal
limiting locations. The same argument applies to neighboring corners
and matched diagonals. Thus all corners of a triangular face receive
uniformly close locations. Join them by short round arcs and fill each
triangle in a small convex ball. This gives $\pi_j$ with vanishing
projected triangle diameter, (T1) from the chronological primal walk,
and (T2) from its visited vertices. For a planar exhaustion, apply
this construction only after supplying a compact sewing that preserves
the tested buffered region.

\begin{lemma}[Uniform topological joining]
\label{bg:small-image-joining}
Under (T1)--(T3), for every $\epsilon>0$ there exist $\delta>0$ and
$j_0$ such that, for $j\ge j_0$ and any $x,y\in\Sigma_j$ with
\[
 d_{\mathrm{sph}}(\pi_j(x),\pi_j(y))<\delta,
\]
there is a path from $x$ to $y$ whose projection is contained in
$B(\pi_j(x),\epsilon)$. If $x,y$ are primal vertices, the path can be
taken in $G_j$. The image of $\pi_j$ is asymptotically dense in
$\mathbb S^2$.
\end{lemma}

\begin{proof}
First take primal vertices. If the assertion fails, there are an
$\epsilon>0$, a subsequence, and vertices
$x_j=c_j(s_j)$, $y_j=c_j(t_j)$ whose projected distance tends to zero
but which cannot be joined as asserted. Pass to limits $s_j\to s$,
$t_j\to t$. By (T1), $\eta(s)=\eta(t)=z$. Choose one finite chain
from (T3). For each of its generators use its discrete approximating
path. Join an approximating endpoint to the endpoint used by the
next generator by the corresponding segment of $c_j$. Those
chronological intervals shrink to the shared limiting time, so their
projected images converge uniformly to $z$, by (T1) and continuity
of $\eta$. Use the same chronological connectors at $s_j,t_j$.
There are only finitely many pieces. Their concatenation eventually
projects into $B(z,\epsilon/2)$, a contradiction.

For arbitrary surface points, first use (T2), then the vertex version
with a smaller error, and concatenate. Finally, (T1) and surjectivity
of $\eta$ give asymptotic density of $\pi_j(\Sigma_j)$.
Notice that chronological connectors in this proof can contain
unboundedly many primal edges.
\end{proof}

\subsection{Connected exteriors, disk fillings, and degree}

The arbitrary-lift joins now let us fill loops without leaving a
prescribed projected neighborhood. The elementary fact behind this
step is that every component of the complement of a nonempty compact
connected set $K$ in a topological sphere is simply connected. Indeed,
a polygonal simple loop disjoint from $K$ has one Jordan disk disjoint
from $K$, because $K$ lies on one side. The Jordan--Schoenflies
Theorem contracts the loop in that disk and hence in its complementary
component. An arbitrary loop in the open complement is homotopic
there to a polygonal loop with finitely many crossings. Split it at
the crossings and contract the resulting simple loops as above.
No local regularity of $K$ is needed.

\begin{lemma}[Exterior connection and small-loop filling]
\label{bg:small-loop-filling}
Suppose continuous maps $\pi_j:\Sigma_j\to\mathbb S^2$ from
topological spheres have asymptotically dense images and satisfy the
uniform joining conclusion of Lemma~\ref{bg:small-image-joining}.
Fix $0<r<R$ and $z\in\mathbb S^2$. For all sufficiently large $j$
there is a nonempty compact connected set $K_j\subset\Sigma_j$
such that
\[
 \pi_j^{-1}(\mathbb S^2\setminus B(z,R))\subset K_j,
 \qquad K_j\cap\pi_j^{-1}(\overline{B(z,r)})=\varnothing.
\]
Consequently every loop projecting into $B(z,r)$ bounds a continuous
disk projecting into $B(z,R)$. The index threshold can be chosen
uniformly in $z$.
\end{lemma}

\begin{proof}
Choose $r<r_1<r_2<R$ and a joining error much smaller than both
$r_1-r$ and $R-r_2$. For large $j$, choose one anchor lift whose
image lies outside $B(z,R)$. Such a lift exists by asymptotic density.
The exterior of $\overline{B(z,r_2)}$ is path connected. Join the
image of any prescribed lift of a point outside $B(z,R)$ to the
anchor image by a target path in that exterior. Subdivide the path
into sufficiently small steps. Asymptotic density supplies lifts
whose images are sufficiently close to the subdivision points;
retain the prescribed lifts at the ends. Apply the uniform joining
property to successive lifts. The resulting upstairs path has
projection outside $\overline{B(z,r_1)}$.

Take the union of the anchor and these paths for every lift in
$\pi_j^{-1}(\mathbb S^2\setminus B(z,R))$, and then take its closure
in $\Sigma_j$. This is a compact connected set $K_j$ containing
that entire preimage. Its projection remains outside $B(z,r_1)$,
so it is disjoint from the preimage of $\overline{B(z,r)}$.
Every loop in the latter preimage lies in a component of
$\Sigma_j\setminus K_j$. That component is simply connected and
is contained in $\pi_j^{-1}(B(z,R))$, proving the filling claim.
The joining and density bounds used here are independent of $z$.
\end{proof}

\begin{proposition}[Degree and robust planar topology]
\label{bg:projection-topology}
Under the hypotheses of Lemma~\ref{bg:small-loop-filling}, the maps
$\pi_j$ have degree $+1$ or $-1$ for all sufficiently large $j$.
More precisely, for every $\epsilon>0$ and sufficiently large $j$
there is a continuous map $g_j:\mathbb S^2\to\Sigma_j$ with
\[
 \sup_{z\in\mathbb S^2}
 d_{\mathrm{sph}}(\pi_j(g_j(z)),z)<\epsilon.
\]
Loops with a nonzero projected winding difference robustly separate
the corresponding lifted fibers, as specified below; in particular,
opposite primal crossings of a rectangle with fixed positive margins
intersect in the primal graph.
\end{proposition}

\begin{proof}[Proof of the approximate right inverse and degree]
Fix a sufficiently small error $\epsilon$. First choose a loop-image
radius much smaller than $\epsilon$, small enough that
Lemma~\ref{bg:small-loop-filling} fills loops of that radius inside
an $\epsilon/3$ ball. Next choose a joining error much smaller than
that loop-image radius, and then a target triangulation so fine that
the endpoints of each of its edges can be joined using
Lemma~\ref{bg:small-image-joining}. For sufficiently large $j$,
choose for each target vertex a lift with projection arbitrarily
close to that vertex. Connect the chosen lifts along each target
edge by the joining property. This defines $g_j$ on the
one-skeleton.

For each triangular face, the image under $\pi_j\circ g_j$ of its
boundary lies in a ball of the chosen loop-image radius about one
of its vertices. Lemma~\ref{bg:small-loop-filling} extends $g_j$
continuously over this face, with projected image inside the
$\epsilon/3$ ball about that vertex. The extensions agree on shared
edges. Fineness of the target mesh gives the asserted uniform error
on every face. Short round geodesics give a homotopy from
$\pi_j\circ g_j$ to the identity. Thus
\[
 \deg(\pi_j)\deg(g_j)=1,
\]
and both integer degrees are $+1$ or both are $-1$. In particular
$\pi_j$ is onto. No embedding or injectivity of $g_j$ is claimed.
\end{proof}

\subsection{Graph intersections forced by winding}

We next use the full strength of joining arbitrary lifts: winding
separates entire fibers, a property that degree alone would not give.
For distinct target points $a,b$, the group
$H_1(\mathbb S^2\setminus\{a,b\};\mathbb Z)$ is infinite cyclic.
A loop has a \emph{nonzero winding difference about $a,b$} if its
class in this group is nonzero. This formulation avoids choosing a
point at infinity.

\begin{lemma}[Separation of entire fibers]
\label{bg:winding-separation}
Under the hypotheses of Lemma~\ref{bg:small-loop-filling},
fix $\rho>0$. For sufficiently large $j$, suppose a loop
$\lambda:\mathbb S^1\to\Sigma_j$ satisfies
\[
 \operatorname{dist}_{\mathrm{sph}}
       (\pi_j(\lambda),\{a,b\})>\rho
\]
and $\pi_j\circ\lambda$ has nonzero winding difference about $a,b$.
Then every path joining any lift of $a$ to any lift of $b$ intersects
$\lambda$. The loop need not be simple.
\end{lemma}

\begin{proof}
All lifts of $a$ can be joined to any one prescribed lift by paths
projecting into $B(a,\rho/3)$, using the uniform joining property.
After closure, their union is compact, connected, contains the
entire fiber of $a$, and is disjoint from $\lambda$. The same holds
for $b$. If an upstairs path connected the prescribed lifts while
avoiding $\lambda$, its union with these two compact connected sets
would be a compact connected set $K$ containing both entire fibers
and disjoint from $\lambda$. The elementary complement fact above
would contract $\lambda$ in $\Sigma_j\setminus K$. Projection of
this contraction would contract $\pi_j\circ\lambda$ in
$\mathbb S^2\setminus\{a,b\}$, contradicting its nonzero homology
class. This proves separation even if exact fibers are disconnected:
their pieces were first joined away from the loop.
\end{proof}

We apply this separation test to transverse graph crossings.
Work in a planar coordinate disk containing $[-3,3]^2$.
Fix $0<a<1/4$. A horizontal passage has projected image in
$[-2,2]\times[-a,a]$ and endpoints on the lines $x=-2,x=2$.
A vertical passage has projected image in
$[-a,a]\times[-2,2]$ and endpoints on the lines $y=-2,y=2$.
The same conclusions hold with small errors in these endpoints and
tubes, by making the margins smaller first. Rectangles obtained by
a fixed change of coordinates are treated identically.

\begin{corollary}[Primal crossing intersection]
\label{bg:primal-crossing}
Assume (T1)--(T3). For sufficiently large $j$, any two primal graph
paths meeting the preceding horizontal and vertical passage tests share
a primal vertex. More generally, the assertion holds for any finite
configuration of tubes and open ports whose margins permit this
test after restricting to subpassages and changing coordinates.
\end{corollary}

\begin{proof}
Call the horizontal path $\alpha$ and the vertical one $\beta$.
Close $\alpha$ upstairs by a path whose projection follows an arc
from its right endpoint to its left endpoint going down outside
$[-a,a]\times[-2,2]$, around the bottom of this rectangle, and up
its other side. Such an upstairs path is obtained by a finite chain
of the small-image joins, starting and ending at the prescribed
lifts. Choose its projected error smaller than the fixed clearance
from the vertical tube. It is consequently disjoint from $\beta$.

Let $b_-$ and $b_+$ be the actual projected endpoints of $\beta$.
The projected closed loop has winding difference one, up to sign,
about $b_-,b_+$. To check this without a simplicity assumption,
homotope the horizontal piece inside its horizontal strip to a
straight left-to-right segment, and homotope the added piece in
its thin exterior tube to the chosen closing arc. These homotopies
avoid both endpoint neighborhoods. The resulting rectangle loop
contains $b_-$ and excludes $b_+$, giving the stated winding
difference. Lemma~\ref{bg:winding-separation} forces $\beta$ to
meet the closed loop. It cannot meet the added piece and therefore
meets $\alpha$. Two paths in an embedded primal graph can intersect
only on a common vertex or a common edge; a common edge also gives
a common vertex. Thus the intersection is an actual primal graph
intersection, at which their walks may be stitched.
\end{proof}

\begin{remark}[Why the arbitrary-lift hypothesis is necessary]
\label{bg:arbitrary-lifts-essential}
Vanishing projected mesh and existence of some nearby lifts do not
suffice. The map $z\mapsto z^2$ on the Riemann sphere has degree two
and admits triangulations with arbitrarily small projected mesh.
Over a disk avoiding its two branch values it has two disjoint
sheets. A horizontal path on one sheet and a vertical path on the
other have robustly crossing projections and are disjoint upstairs.
Nearby target points have suitably chosen nearby lifts, but two
prescribed lifts of the same point cannot be joined with projection
confined to that disk. Lemma~\ref{bg:small-image-joining} rules out
exactly this obstruction.
Asymptotic density of the image is also necessary: a constant map
satisfies the small-image joining property but has degree zero.
\end{remark}

\subsection{Cellular neighborhoods and retained stars}

The crossing test gives graph intersections. We also need a disk
around a prescribed fiber whose nearby primal and dual incidences
remain in retained time blocks. A disk preimage, or the union of its
incident triangles, may have holes and pinches. We first remove
these inside a larger collar, then control the necessary cell stars.

For a subset $A\subset\Sigma_j$, let $\mathcal S_j^0(A)$ be the union
of all closed original triangles meeting $A$, with their faces. Let
$\mathcal S_j^{r+1}(A)$ be the union of the closed original triangles
sharing a vertex with a triangle of $\mathcal S_j^r(A)$. Thus only a
specified finite number of full cell-star layers is retained; a proper patch is not required to contain stars of all its
vertices indefinitely. Write $\mu_j$ for the maximum projected
diameter of an original triangle. For each fixed $r$,
\begin{equation}\label{bg:star-image-bound}
 \pi_j(\mathcal S_j^r(A))
 \subset \{x:d_{\mathrm{sph}}(x,\pi_j(A))\le(r+1)\mu_j\}.
\end{equation}
Indeed, successive triangles share the image of a vertex, so their
diameter bounds can be added. Vertex degrees need not be bounded.

\begin{lemma}[A cellular disk between two lifted disks]
\label{bg:cellular-disk}
Assume (T1)--(T3) and $\mu_j\to0$. If $E\Subset\operatorname{int}F$
are closed topological disks in $\mathbb S^2$, then, for sufficiently
large $j$, there is a closed topological disk $\Delta_j\subset\Sigma_j$
such that
\[
 \pi_j^{-1}(E)\subset\operatorname{int}\Delta_j,
 \qquad
 \Delta_j\subset\pi_j^{-1}(\operatorname{int}F).
\]
The disk can be taken to be a subcomplex of a finite subdivision
of the original triangulation. This subdivision supplies a topological
boundary only; it adds no edges to the primal metric graph.
\end{lemma}

\begin{proof}
Choose nested disk collars strictly between $E$ and $F$. We connect
the lifted core, take a regular neighborhood, and fill its holes
except for the component containing the prescribed exterior.

\emph{Connect the core.} For large $j$, the projection is onto by
Proposition~\ref{bg:projection-topology}. Connect every lift over $E$
to one anchor by finite chains of small-image joins along target
paths in $E$. Choose the joining error smaller than the first collar.
The closure of the union is a compact connected set $H_j$ containing
$\pi_j^{-1}(E)$ and projecting inside that collar. The triangles
meeting $H_j$ form a connected finite subcomplex, whose full cell
star lies in the next collar by Equation~\ref{bg:star-image-bound}.

\emph{Take a regular neighborhood.} Barycentrically subdivide the
triangulation. Give a new vertex value zero if its original simplex
belongs to the connected subcomplex, and value one otherwise.
Extend linearly and let $N_j$ be the sublevel set at $1/2$.
The boundary avoids new vertices and cuts each crossed triangle in
a segment, with exactly two segments meeting across each shared edge.
It is therefore a disjoint union of polygonal circles, and $N_j$ is
a compact subsurface. Normalizing the barycentric weights of the
zero-valued vertices retracts each piece to its zero face. These
retractions agree on common faces, so $N_j$ is connected and contains
the original subcomplex in its interior. It lies in the full original
cell star. Subdivide along the level segments and triangulate the
resulting polygons to make $N_j$ cellular. This resolves pinches at
original vertices without modifying the primal graph.

\emph{Fill the interior holes.} The exterior-connection construction
of Lemma~\ref{bg:small-loop-filling}, using the intermediate disk
collars in place of round disks, gives a compact connected set $K_j$
containing
$\pi_j^{-1}(\mathbb S^2\setminus\operatorname{int}F)$ and disjoint
from $N_j$. The proof is unchanged: the exterior of a topological
disk is path connected, and each of the finitely many collars has
positive separation from the next. All of $K_j$ lies in one component
$O_j$ of $\Sigma_j\setminus N_j$.

A connected proper compact subsurface of a sphere has complementary
components which are open disks, one for each boundary component.
This follows directly by applying the Jordan--Schoenflies Theorem to
its disjoint boundary circles and using connectedness of the
subsurface. Fill every complementary component of $N_j$ except
$O_j$, and set $\Delta_j=\Sigma_j\setminus O_j$. This is a closed
disk containing $H_j$ in its interior. It avoids $K_j$, hence projects
inside $\operatorname{int}F$. Filling the complementary components
preserves cellularity after the same subdivision.
\end{proof}

The cellular disk can now be placed inside retained time blocks.
For each original triangle, let $I_{j,\sigma}$ be its chronological
mesh interval. We use the following projection convention:
\begin{equation}\label{bg:triangle-time-control}
 \max_{\sigma}\ \sup_{x\in\sigma,\ t\in I_{j,\sigma}}
 d_{\mathrm{sph}}(\pi_j(x),\eta(t))\longrightarrow0,
 \qquad
 \max_{\sigma}|I_{j,\sigma}|\longrightarrow0.
\end{equation}
The continuous projection construction above has this property.
Subdivided cells inherit the interval of their original triangle.

\begin{proposition}[Cellular sewing stability with margins]
\label{bg:sewing-stability}
Suppose admissible discrete contour excursions have actual sphere
cell complexes, converge uniformly to contours whose finite-chain
quotient is a sphere, and satisfy the link approximation and
chronological conditions (T1)--(T3) and
Equation~\ref{bg:triangle-time-control}. In particular, the scalar
approximation hypothesis includes any separately prescribed generators
at exceptional or distinguished times.

Let $z\in\mathbb S^2$ have finite fiber $S=\eta^{-1}(z)$, and let $J$
be a finite union of open retained time intervals containing $S$.
Choose closed topological disks about $z$ with
\[
 E_{-1}\Subset\operatorname{int}E_0\Subset\operatorname{int}E_1
 \Subset\operatorname{int}E_2\Subset\operatorname{int}E_3
 \Subset\operatorname{int}E_4,
 \qquad \eta^{-1}(E_4)\subset J.
\]
For every fixed integer $r\ge0$ and sufficiently large $j$, there
are nested cellular disks $\Delta_j^-\Subset\operatorname{int}\Delta_j^+$
such that
\begin{align*}
 \pi_j^{-1}(E_0)&\subset\operatorname{int}\Delta_j^-
       \subset\Delta_j^-\subset\pi_j^{-1}(\operatorname{int}E_1),\\
 \pi_j^{-1}(E_2)&\subset\operatorname{int}\Delta_j^+
       \subset\Delta_j^+\subset\pi_j^{-1}(\operatorname{int}E_3).
\end{align*}
Writing $A_j=\Delta_j^+\setminus\operatorname{int}\Delta_j^-$:
\begin{enumerate}[label=(\roman*)]
\item $A_j$ is a closed cellular annulus and contains
$\pi_j^{-1}(E_2\setminus\operatorname{int}E_1)$.
\item Every original triangle in $\mathcal S_j^r(\Delta_j^+)$ has
its entire chronological interval in $J$. In particular, taking
$r\ge1$ retains complete primal and dual vertex stars of every
original triangle meeting the disk.
\item Every original triangle in $\mathcal S_j^r(A_j)$ has its
chronological interval in a fixed compact subset of $J\setminus S$,
up to an arbitrarily small fixed open time enlargement still contained
there. It can therefore be covered by finitely many off-center time
blocks, chosen with positive margins.
\item Robust transverse primal crossings in the annular region have
the actual intersections asserted by Corollary~\ref{bg:primal-crossing}.
A closed primal walk there with nonzero projected winding difference
about $z$ and a point of $\mathbb S^2\setminus E_4$ contains a simple
primal circuit surrounding $\pi_j^{-1}(E_0)$. Its chosen interior
is contained in $\Delta_j^+$,
so its interior and the required fixed number of cell-star layers
lie in the retained patch.
\end{enumerate}
All conclusions are topological. They bound neither the length of an
auxiliary joining path nor a normalized endpoint cost.
\end{proposition}

\begin{proof}
We first construct the annulus, then retain its incidence layers and
extract a separating primal circuit.
Apply Lemma~\ref{bg:cellular-disk} to $(E_0,E_1)$ and $(E_2,E_3)$.
Since $E_1\subset\operatorname{int}E_2$, the resulting disks are
nested as claimed. Their difference is an annulus; the displayed
inclusions give (i).

For (ii), Equation~\ref{bg:star-image-bound} and the margin between
$E_3,E_4$ place all the star triangles, and their chronological
images in Equation~\ref{bg:triangle-time-control}, inside $E_4$ for
large $j$. The compact set $\eta^{-1}(E_4)$ is contained in the open
set $J$. Vanishing chronological mesh therefore places each whole
triangle interval in $J$, not just one representative time.

Every point of $A_j$ projects outside $E_0$ and inside $E_3$.
The margins to $E_{-1}$ and $E_4$ show that every triangle in its
fixed star buffer has chronological images in
$E_4\setminus\operatorname{int}E_{-1}$ for large $j$. The inverse
image of this compact set is a compact subset of $J$ disjoint from
$S$. A slightly larger compact time neighborhood still in
$J\setminus S$ contains the full triangle intervals. A finite
interval cover, with slightly larger buffers still off $S$, proves
(iii). Any additional finite number of incidence guards is covered
by increasing $r$ before taking $j$ large.

For (iv), use Corollary~\ref{bg:primal-crossing} for transverse
crossings. To obtain the surrounding circuit, decompose a closed primal
walk into finitely many simple circuits. Projected winding is additive,
so at least one has nonzero winding difference. By
Lemma~\ref{bg:winding-separation}, this circuit separates a lift
of $z$ from a lift of an exterior point. It lies in $A_j$, whereas
$\operatorname{int}\Delta_j^-$ is connected and disjoint from it;
therefore all of $\pi_j^{-1}(E_0)$ lies on its inner side. The
complement of $\Delta_j^+$ is connected, lies on the exterior side,
and contains the exterior lift. The chosen circuit interior is
thus contained in $\Delta_j^+$. Its fixed star buffers are retained
by (ii). The circuit uses only edges of the original walk.
\end{proof}

\subsection{Capping while preserving local incidence}

The preceding arguments apply to compact sphere sewings. To use them
in a planar exhaustion, we must close a bounded contour window without
changing its protected incidences. We first build admissible discrete
caps. For the limiting quotient, we then use Brownian caps whose law
is absolutely continuous with respect to an excursion law.

\begin{lemma}[Axial completion as an actual sphere map]
\label{bg:axial-completion}
A finite walk in $\mathbb Z_{\ge0}^2$ with steps
$\{(1,0),(-1,0),(0,1),(0,-1)\}$, starting and ending at $(0,0)$,
defines an actual connected primal tree/cotree map on an oriented
sphere. A down-step is matched to the most recent unmatched up-step
of its own coordinate. The associated triangle and quadrangle cell
complex, rather than the literal quotient of the chronological
circle, is the discrete surface.

For lattice endpoints $a,b$, an axial filler of $m$ steps is possible
only if
\[
 m\ge \lVert b-a\rVert_1,
 \qquad m\equiv(b_1-a_1)+(b_2-a_2)\pmod 2.
\]
These conditions are sufficient with a nonnegative filler when both
endpoints are nonnegative. Moreover, a continuous filler which is
positive on its open time interval can be uniformly approximated,
under height scale $\sqrt n$ and time scale $n$, by such fillers of
prescribed positive macroscopic duration and compatible rounded
endpoints. Strict positivity holds on every fixed interior time
compact for all sufficiently large $n$.
\end{lemma}

\begin{proof}
\emph{The sphere map.} The first-coordinate up-steps and their stack
matches give a rooted plane tree. Thicken it on the sphere. Its
complement is a disk following the tree contour, with corners marked
by the intervening second-coordinate steps. Draw an arc in this disk
for each matched second-coordinate pair. Stack matching makes the
arcs noncrossing; their endpoints are the specified primal corners,
with loops and repeated endpoints allowed. Cutting along the arcs
leaves disks, giving a cellular sphere embedding. The arcs are the
primal diagonals of the sewn quadrangles, and their duals form the
complementary dual tree. If there are $n_1,n_2$ up-steps of the two colors, respectively, the map
has $n_1+1$ vertices, $n_1+n_2$ edges, and $n_2+1$ faces, consistently
with this construction. The case of a one-vertex initial tree is
included. Notice that the word with steps $(1,0),(-1,0)$ still gives
a sphere map; its identically zero second contour would collapse
the literal time-circle quotient to a point.

\emph{Discrete fillers.} First move each coordinate monotonically from $a$
to $b$. This takes $\lVert b-a\rVert_1$ steps and stays nonnegative.
Any even excess is supplied by cancelling positive-then-negative
pairs in one coordinate. This proves sufficiency and necessity.

\emph{Approximation of continuous fillers.} Fix a polygonal time
partition on which the continuous filler oscillates by less than
$\epsilon$. Round the positions at its partition points to the
required parity lattice.
On each block the available number of steps is of order $n$, whereas
the endpoint displacement is of order $\sqrt n$. The preceding
inequality and parity conditions therefore hold for large $n$.
Use monotone coordinate moves between the rounded endpoints and
cancelling pairs for the excess. The path stays in their coordinate
rectangle, enlarged by one height-mesh unit. This is within a fixed
multiple of $\epsilon$ of the continuous path throughout the block.
Interior positivity follows from its positive minimum on compact
interior time intervals. At an excursion endpoint the first two
steps, or the last two steps, may be arranged to enter, or leave,
the strictly positive quadrant; these modifications have vanishing
time and height size. Let $\epsilon\downarrow0$ by a diagonal
choice of partitions. Even cuts and parity-preserving translations
give the required compatibility at retained block boundaries.
\end{proof}

\begin{lemma}[Protected capping of a bounded window]
\label{bg:protected-capping}
Consider a two-sided resolved contour walk and a finite set of its
triangles. Fix a finite number of cell-star layers to be protected.
For each coordinate, suppose its values are unbounded below on both
ends of the time axis. Then one can retain a larger
bounded window and complete it by nonnegative axial prefixes and
suffixes to an admissible finite excursion so that the protected
triangles and their prescribed full-star layers have exactly the
same primal, dual, side, and corner incidences. The prefix and suffix
introduce no new partner for a protected step or corner.

At the continuum level, a compact increment block whose law is
absolutely continuous with respect to the corresponding Brownian
increment law can be completed, up to a common translation of each
coordinate, to a positive quadrant excursion whose whole law is
absolutely continuous with respect to a longer Brownian quadrant
excursion. Thus a compact capping used for the projection argument
can have a sphere quotient. Arbitrary deterministic positive fillers
alone do not imply that quotient conclusion.
\end{lemma}

\begin{proof}
\emph{Protect the discrete incidences.} Start with a bounded window
containing the desired triangles. In each coordinate find a visit to
the left and a visit to the right strictly below every corner height and adjacent step level used in
that window. Retain a window beyond all four visits. The two colors
need not attain their low values at the same times. A scalar chord
or matched step from the original window to outside the enlarged
one would cross its coordinate's lower visit, which is impossible.
Consequently every scalar partner needed for one full incidence
layer is now in the enlarged window. Repeat this finite enlargement
for the prescribed number of layers, and once more to protect their
incidences. In a contour coupling, take strictly positive macroscopic
height margins at these lower visits; uniform convergence preserves
the inequalities, including adjacent lattice levels.

Keep this entire final window unchanged. Translate each coordinate
by a sufficiently large lattice constant so that the window lies
in the positive quadrant, respecting the parity lattice. Attach
nonnegative axial paths from zero to its initial point and from its
final point back to zero. Lemma~\ref{bg:axial-completion} gives an
actual sphere map. All relations inside the retained window use the
same equalities and intervening minima, and the retained lower
visits exclude any partner in a cap. This proves the exact incidence
claim. Independent choices of lower visits in the two coordinates
suffice throughout.

\emph{Preserve the continuum sphere quotient.} Let $W$ be an
increment block of duration $T$, and consider a Brownian
quadrant excursion of duration $T+2$. The law of its middle increment
block, from time $1$ to time $T+1$, has a strictly positive density
with respect to Brownian increments of duration $T$. Indeed, the
Markov bridge decomposition writes that density, up to normalization,
as
\[
 g(w)=\int_{\mathbb R_{>0}^2}
 e_1(a)e_1(a+w(T))
 \boldsymbol 1_{\{a+w(t)\in\mathbb R_{>0}^2\ \text{for all }0\le t\le T\}}
 \,da,
\]
where $e_1$ is the positive entrance density, equal to the exit
density by reversal, with the fixed Brownian covariance understood.
These densities have Gaussian decay times a fixed harmonic power.
The integral is finite and is positive for every bounded continuous
$w$, since all sufficiently large coordinate translations place
its range strictly inside the quadrant. This is also the usual
killed Brownian bridge decomposition: the indicator supplies the
killing during the middle block.

Sample the start height and the two caps from this excursion law
conditionally on its middle increments being $W$. Since the given
law of $W$ is absolutely continuous with respect to Brownian
increments, and $g>0$, the resulting entire excursion law is
absolutely continuous with respect to the reference excursion law.
Its Radon--Nikodym derivative is the density of the prescribed middle
increments divided by the normalized middle-increment density of the
reference excursion. The middle block is reproduced exactly up to
translation. Its caps are continuous and positive in their interiors.
Theorem~\ref{bg:continuum-inputs}(ii)--(iii) gives the sphere quotient, and
Lemma~\ref{bg:axial-completion} approximates these caps by admissible
axial fillers while retaining the actual middle steps. A countable
choice among deterministic bounded windows does not change the
absolute-continuity assertion, after conditioning on each chosen
window of positive probability.
\end{proof}

We finally identify the protected neighborhood before and after
capping. Suppose a neighborhood in the original proper planar
quotient has its \emph{entire} chronological preimage in the initial
protected time window, inside all retained lower barriers. Its
topology agrees with that of the corresponding capped neighborhood.
Every first scalar link from such a fiber is preserved and has no
new exterior partner; induction along finite chains therefore gives
the same complete fibers. A saturated open chronological set has
open image under either quotient map, and the two restrictions are
quotients by the same equivalence relation. They are consequently
homeomorphic. This is the local identification used in an exhaustion;
it does not identify arbitrary unprotected cap geometry with the
original plane. For distinguished auxiliary profiles whose laws
are singular at prescribed times, use their separately proved Moore
quotient criterion and generator approximations in
Proposition~\ref{bg:sewing-stability}; the Brownian absolute-continuity
completion is not asserted for those singular profiles.
